\documentclass[12pt,a4paper]{report}

\usepackage[utf8]{inputenc}
\usepackage[T1]{fontenc}
\usepackage[serbian,english]{babel}
\usepackage[a4paper,margin=2.5cm]{geometry}
\usepackage{setspace}
\usepackage{hyperref}
\usepackage{longtable}
\usepackage{booktabs}

\onehalfspacing

\begin{document}

\begin{titlepage}
\begin{center}
UNIVERZITET U NOVOM SADU \\
FAKULTET TEHNIČKIH NAUKA U NOVOM SADU

\vfill

{\Large [Ime i prezime kandidata]}

\vfill

{\Large \textbf{Razvoj 2D gejm endžina sa editorom scena i optimizovanom detekcijom kolizija zasnovanom na segmentisanoj sweep-and-prune metodi}}

\vfill

Diplomski rad \\
Osnovne akademske studije

\vfill

Novi Sad, [2026.]
\end{center}
\end{titlepage}

\chapter*{Zadatak rada}

Student: [Ime i prezime] \\
Broj indeksa: [SW xx/20xx] \\
Mentor: [Ime i prezime, zvanje]

\begin{enumerate}
  \item Analizirati stanje u oblasti 2D gejm endžina, editora scena i broad-phase detekcije kolizija.
  \item Definisati arhitekturu sistema koja obuhvata runtime jezgro, fizički podsistem i editor scena.
  \item Implementirati softversko rešenje i demonstracioni primer igre.
  \item Izvršiti funkcionalnu i, po mogućnosti, kvantitativnu evaluaciju sistema.
  \item Dokumentovati rezultate, ograničenja i pravce daljeg razvoja.
\end{enumerate}

\tableofcontents

\chapter{Uvod}

U ovom poglavlju treba predstaviti problem rada, motivaciju, ciljeve, doprinose i strukturu rada. Kao osnovu je moguće iskoristiti tekst iz fajla \texttt{paper/UVOD\_FORMALNI\_PLATFORMATOR.md}.

\chapter{Pregled stanja u oblasti}

U ovom poglavlju treba obraditi:

\begin{itemize}
  \item broad-phase algoritme za detekciju kolizija,
  \item sweep-and-prune i inkrementalne pristupe,
  \item segmentisane intervalne liste iz rada Tracy, Buss i Woods,
  \item komponentnu arhitekturu u savremenim gejm endžinima,
  \item Unity i Godot kao referentne alate za tok rada sa scenama.
\end{itemize}

\chapter{Teorijske osnove}

U ovom poglavlju treba opisati komponentnu arhitekturu, AABB omotače, grid-subdiviziju, segmentisane intervalne liste, SAT pristup i impulsno rešavanje kontakata.

\chapter{Projektovanje i implementacija sistema}

U ovom poglavlju treba obraditi sledeće celine:

\begin{itemize}
  \item arhitekturu runtime jezgra,
  \item GameManager, GameObject i sistem komponenti,
  \item PhysicsManager, grid i AABB strukture,
  \item segmentisane intervalne liste i checkpoint mehanizam,
  \item SAT i impulsni solver,
  \item editor scena i njegov korisnički interfejs,
  \item Mario primer kao potrošača biblioteke.
\end{itemize}

\chapter{Evaluacija i rezultati}

Ovo poglavlje treba da sadrži opis test scenarija, funkcionalnu validaciju, pregled izvršenih projektnih testova, diskusiju rezultata i posebnu tabelu sa benchmark merenjima broad-phase podsistema kada ta merenja budu dostupna.

\chapter{Demonstracija sistema}

Ovde treba opisati tok demonstracije editora, runtime scene i Mario primera, sa fokusom na ono što komisija može direktno da vidi tokom odbrane.

\chapter{Zaključak}

Zaključak treba da sumira ostvarene ciljeve, doprinos rada, ograničenja implementacije i pravce budućeg razvoja.

\begin{thebibliography}{10}

\bibitem{tracy2009}
Tracy, D. J., Buss, S. R., Woods, B. M.
Efficient Large-Scale Sweep and Prune Methods with AABB Insertion and Removal.
IEEE Virtual Reality Conference, 2009.

\bibitem{baraff1992}
Baraff, D.
Dynamic Simulation of Non-Penetrating Rigid Bodies.
Cornell University, 1992.

\bibitem{cohen1995}
Cohen, J. D., Lin, M. C., Manocha, D., Ponamgi, M. K.
I-COLLIDE: An Interactive and Exact Collision Detection System for Large-Scale Environments.
Symposium on Interactive 3D Graphics, 1995.

\bibitem{bergen2003}
van den Bergen, G.
Collision Detection in Interactive 3D Environments.
Morgan Kaufmann, 2003.

\bibitem{unity}
Unity Technologies.
Unity Manual.
\url{https://docs.unity3d.com/Manual/}

\bibitem{godot}
Godot Engine.
Official Documentation.
\url{https://docs.godotengine.org/}

\end{thebibliography}

\chapter*{Biografija}

[Ovde uneti biografiju kandidata u skladu sa FTN smernicama.]

\chapter*{Ključna dokumentacijska informacija}

Autor: [Ime i prezime kandidata] \\
Mentor: [Ime i prezime mentora, zvanje] \\
Naslov rada: Razvoj 2D gejm endžina sa editorom scena i optimizovanom detekcijom kolizija zasnovanom na segmentisanoj sweep-and-prune metodi \\
Ključne reči: gejm endžin, detekcija kolizija, sweep and prune, editor scena, komponentna arhitektura

\chapter*{Key words documentation}

Author: [Candidate full name] \\
Mentor: [Mentor full name, title] \\
Title: Development of a 2D Game Engine with a Scene Editor and Optimized Collision Detection Based on the Segmented Sweep-and-Prune Method \\
Keywords: game engine, collision detection, sweep and prune, scene editor, component-based architecture

\end{document}